%% kapa: exam=2 round=31 year=2020 period=1 subject=appraisal_practice subject_ko=감정평가실무 seq=2 number=2 points=30
%% source: Q-Net 감정평가사 기출문제 (cst00309) artlSeq=5207673 file=raw/1008/5207673/1교시_감정평가실무.pdf
%% conversion: kordoc 4.24.51 -> tools/gen2.py (auto); review=proofread-gpt-astra; flags=none
\begin{longquestion}{【문제 2】}
감정평가사甲은 A자산운용사로부터 부동산 펀드에 새로이 편입되는 Z마트(대형할인점) 3개 점포의 금융기관 담보제공 목적 감정평가를 의뢰받았다. 관련 법규 및 이론을 참작하고, 제시된 자료 및 전제조건을 활용하여 다음의 물음에 답하시오.(단, 기준시점은 2020. 09. 19임) (30점) \par
\medskip\noindent 물음 1) 대상물건 각 점포에 적용되어야 하는 실무적ㆍ이론적으로 타당한 할인율과 재매도환원율을 결정하고, 그 사유를 서술하시오.(단, 할인율과 재매도환원율은 백분율 기준 소수점 둘째자리에서 반올림함) (15점) \par
\medskip\noindent 물음 2) 결정된 할인율과 재매도환원율을 적용하여 대상물건의 수익환원법 시산가액을 산정한 후 원가법에 의한 시산가액과 비교ㆍ검토하고, 각 점포별 시산가액의 균형에 대해 서술하시오. (15점) \par
< 자료 1 > 대상물건 내역 및 원가법 시산가액 \par
{\small\centering\renewcommand{\arraystretch}{1.25}
\resizebox{\ifdim\width>\linewidth\linewidth\else\width\fi}{!}{
\begin{tabular}{|c|c|c|c|c|c|}
\hline
\makecell{해당\\자산} & \makecell{소재지역} & \makecell{규모} & \makecell{용도지역} & \makecell{사용승인일} & \makecell{원가법 시산가액\\(단위: 백만원)} \\
\hline
\makecell{Z마트\\a점포} & \makecell{A} & \makecell{대지면적: 3,866m\textsuperscript{2}\\연면적: 18,500m\textsuperscript{2}} & \makecell{준주거} & \makecell{2006.05.31} & \makecell{토지: 57,990 (79.8\%)\\건물: 14,652 (20.2\%)\\합계: 72,642 (100.0\%)} \\
\hline
\makecell{Z마트\\b점포} & \makecell{B} & \makecell{대지면적: 6,520m\textsuperscript{2}\\연면적: 28,000m\textsuperscript{2}} & \makecell{유통상업} & \makecell{2011.12.10} & \makecell{토지: 37,816 (59.4\%)\\건물: 25,872 (40.6\%)\\합계: 63,688 (100.0\%)} \\
\hline
\makecell{Z마트\\c점포} & \makecell{C} & \makecell{대지면적: 12,630m\textsuperscript{2}\\연면적: 43,000m\textsuperscript{2}} & \makecell{유통상업} & \makecell{2015.11.15} & \makecell{토지: 17,682 (30.8\%)\\건물: 39,732 (69.2\%)\\합계: 57,414 (100.0\%)} \\
\hline
\end{tabular}}\par}

※ 각 점포의 건물은 관련 법규에서 정하여진 허용 용적률을 전부 사용하여 건축 되었음. 원가법 시산가액은 「대ㆍ중소기업 상생협력 촉진에 관한 법률」에 의한 사업조정비용 등 무형적 비용을 제외한 금액임 \par
< 자료 2 > 전제조건 등 \par
1. 대형할인점 운영기업은 유동성 확보 차원에서 최근 10년간 주요 대형할인점 점포를 매각후 재임차(Sales and Lease Back) 하였으며, 매각후 재임차 된 점포는 수익환원법 가치를 바탕으로 자산운용사에 매매되는 관행이 성립되었음. 자산운용사 및 기타 시장참여자(재무적 투자자 등)는 위험회피자임 \par
2. 해당 펀드는 5년 후 청산을 목적으로 하는 펀드로서, 펀드에 대해 투자하는 금융기관은 펀드 만기에 대상물건 각 점포별 재매도 가치를 중요하게 생각하는 바, 의뢰인 A자산운용사는 5년 후 재매도 가치가 할인현금수지에 명확히 포함되는 평가모형을 사용하여 줄 것을 감정평가사甲에게 요청하였음 \par
3. 감정평가사甲은 시장관행에 따라 순영업소득(NOI)을 기초로 할인현금수지법(DCF Method)을 적용하기로 하되, 재매도 가치는 내부추계법을 사용하기로 하였음 \par
4. 감정평가사甲이 소속된 D감정평가법인의 심사위원회는 감정평가사甲에게 “Z마트의 경우 재무적 상황이 악화되어 임차계약 연장이 불투명하고, 의뢰인의 목표수익률도 자산별로 상이하므로, 해당 자산이 소속된 지역의 부동산 상황ㆍ자산별 특성 및 판매시설 운용추이 등을 종합적으로 검토하여 자산별로 할인율과 재매도환원율을 달리 사용할 것”을 권고하였음 \par
5. 해당 부동산의 임대차 내역: 대상 물건은 Z마트 a, b, c 점포로서, 다음은 매도인 Z마트(임차인)와 매수인 A자산운용사가 설립한 특수목적법인 A사모부동산투자신탁 제1호(임대인) 간에 체결된 임대차 내역을 요약한 표임 \par
{\small\centering\renewcommand{\arraystretch}{3.5}
\resizebox{\ifdim\width>\linewidth\linewidth\else\width\fi}{!}{
\begin{tabular}{|c|c|r|p{0.51\linewidth}|}
\hline
\makecell{해당\\자산} & 소재지역 & 연간 임대료(원) & \centering 비고\arraybackslash \\
\hline
\makecell{Z마트\\a점포} & A & 3,800,000,000 & \multirow{3}{=}{임대 기간: 2020.09.19～2030.09.18\par 기타 임대 조건:\par ① 대상 물건에 대한 운영경비는 임차인이 전부 부담하는 순임대차임\par ② 연간 임대료는 1년 단위 임대기간 말 후취 조건임\par ③ 연간 임대료는 매년 1.5\%씩 인상함} \\
\cline{1-3}
\makecell{Z마트\\b점포} & B & 3,600,000,000 & \\
\cline{1-3}
\makecell{Z마트\\c점포} & C & 3,500,000,000 & \\
\hline
\end{tabular}}\par}

6. 보증금과 임대차기간 동안의 공실 및 기타수입은 없는 것으로 간주하며, 연간 임대료는 순영업소득과 일치함 \par
7. 할인율 및 재매도환원율을 공공기관 통계를 기초로 결정하는 경우 최근의 경제상황을 고려하여 최근 1년 평균치와 최근 5년 평균치의 중앙값을 순영업소득 및 재매도 가치에 적용하며, 자산운용사 제시 자료를 기초로 할인율 및 재매도환원율을 결정하는 경우 타인자본 차입비율(L/V: 65\%)을 고려한 가중평균수익률(종합할인율)을 순영업소득 및 재매도 가치에 적용함 \par
8. 해당 부동산 펀드는 유보 없이 배당가능금액 전부를 배당하며, 편의상 부동산 영업경비를 제외한 펀드 운용비용 등은 없는 것을 전제로 함 \par
< 자료 3 > 취득 자료 및 조사사항 \par
1. 공공기관에서 발표한 대형 상업용 부동산 수익률 통계는 다음과 같음. 소득수익률 통계는 각 지역에서 여러 개의 표본에 대해 최근 1년 및 5년의 원본가치 대비 순영업소득 비율을 취합한 결과로서, 원본가치는 토지 및 건물을 각각 산정하여 합산하였음 \par
{\small\centering\renewcommand{\arraystretch}{1.25}
\resizebox{\ifdim\width>\linewidth\linewidth\else\width\fi}{!}{
\begin{tabular}{|c|c|c|c|}
\hline
\makecell{구분} & \makecell{소득수익률\\(최근 1년 평균)} & \makecell{투자수익률\\(최근 1년 평균)} & \makecell{자본이득률\\(최근 1년 평균)} \\
\hline
\makecell{A지역} & \makecell{5.00\%} & \makecell{6.50\%} & \makecell{1.50\%} \\
\hline
\makecell{B지역} & \makecell{4.80\%} & \makecell{5.80\%} & \makecell{1.00\%} \\
\hline
\makecell{C지역} & \makecell{4.70\%} & \makecell{5.20\%} & \makecell{0.50\%} \\
\hline
\end{tabular}}\par}

{\small\centering\renewcommand{\arraystretch}{1.25}
\resizebox{\ifdim\width>\linewidth\linewidth\else\width\fi}{!}{
\begin{tabular}{|c|c|c|c|}
\hline
\makecell{구분} & \makecell{소득수익률\\(최근 5년 평균)} & \makecell{투자수익률\\(최근 5년 평균)} & \makecell{자본이득률\\(최근 5년 평균)} \\
\hline
\makecell{A지역} & \makecell{6.20\%} & \makecell{7.70\%} & \makecell{1.50\%} \\
\hline
\makecell{B지역} & \makecell{6.00\%} & \makecell{7.10\%} & \makecell{1.10\%} \\
\hline
\makecell{C지역} & \makecell{5.60\%} & \makecell{6.00\%} & \makecell{0.40\%} \\
\hline
\end{tabular}}\par}

2. A자산운용사가 제시한 해당 펀드의 목표 배당수익률 및 금융기관 대출금리는 다음 표와 같음. 대출금리는 5년간 고정금리이며, 펀드설정기간 동안 원금상환은 없음. 감정평가사甲이 수집한 자료를 통해 검증해 본 결과 해당 자산운용사의 목표수익률은 유사지역 동종 부동산 펀드 목표수익률과 유사한 수준이며, 실현가능성이 매우 높은 것으로 판단됨 \par
{\small\centering\renewcommand{\arraystretch}{1.25}
\resizebox{\ifdim\width>\linewidth\linewidth\else\width\fi}{!}{
\begin{tabular}{|c|c|c|c|}
\hline
\makecell{구분} & \makecell{초기(초년도)\\목표배당수익률\\(매각차익 배당 제외)} & \makecell{장기 목표배당수익률\\(매각차익 배당 포함)} & \makecell{금융기관 대출금리} \\
\hline
\makecell{Z마트\\a점포} & \makecell{6.80\%} & \makecell{8.00\%} & \makecell{3.50\%} \\
\hline
\makecell{Z마트\\b점포} & \makecell{7.50\%} & \makecell{8.30\%} & \makecell{3.50\%} \\
\hline
\makecell{Z마트\\c점포} & \makecell{7.80\%} & \makecell{8.60\%} & \makecell{3.50\%} \\
\hline
\end{tabular}}\par}

3. 감정평가사甲이 분석한 해당 점포 소재지역의 지역분석 내용은 다음과 같음 \par
{\small\centering\renewcommand{\arraystretch}{1.25}
\resizebox{\ifdim\width>\linewidth\linewidth\else\width\fi}{!}{
\begin{tabular}{|c|p{0.83\linewidth}|}
\hline
\makecell{구분} & 지역분석 내용 \\
\hline
\makecell{A지역} & A지역은 S시의 남쪽에 위치하며, 해당 지역의 배후지는 S시 내에서 상대적으로 고소득 계층이 거주하는 지역으로서, 판매시설의 매출은 상대적으로 견고한 추이를 보이고 있음. A지역의 상업용 부동산 임대료는 과거 타지역에 비해 상대적으로 낮은 변동성을 보이고 있으며, 장래 이자율ㆍ환율․GDP 상승률 등 거시경제지표 변동과 상관계수가 낮을 것으로 예상됨 \\
\hline
\makecell{B지역} & B지역은 S시의 서남쪽에 위치하며, 해당 지역의 배후지는 S시 내에서 상대적으로 저소득 계층이 거주하는 지역으로서, 판매시설의 매출은 인근 경공업 경기에 비교적 민감하게 반응하고 있음. B지역의 과거 상업용 부동산 임대료의 변동성은 A지역에 비해 상대적으로 높고 C지역에 비해 상대적으로 낮으며, 장래 B지역 임대료의 이자율ㆍ환율․GDP 상승률 등 거시경제지표와의 상관계수도 A지역에 비해 상대적으로 높고 C지역에 비해 상대적으로 낮을 것으로 예상됨 \\
\hline
\makecell{C지역} & C지역은 K도 남부에 위치하며, 해당 지역의 배후지는 과거 K도 최고수준 소득 계층이 거주하는 지역이었으나, 최근 조선업의 불황으로 인구변동이 활발히 일어나고 있음. C지역 판매시설의 매출은 인근 조선업 경기에 매우 민감하게 반응하고 있음. C지역 상업용 부동산 임대료는 과거 B지역에 비해 상대적으로 높은 변동성을 보이고 있으며, 장래 이자율ㆍ환율․GDP 상승률 등 거시경제지표 변동과 상관계수도 B지역에 비해 상대적으로 높을 것으로 예상됨 \\
\hline
\end{tabular}}\par}

< 자료 4 > 기타사항 \par
1. 할인현금흐름수지표는 십만원 단위에서 반올림하여, 백만원 단위로 작성할 것 \par
2. 재매도 비용은 매각자문비용 등으로서, 재매도 가치의 1.3\%를 적용함 \par
3. 현시점의 기입환원율(Going-in Cap-rate)과 펀드 자산 매각시기의 재매도환원율은 동일한 것으로 간주할 것 \par
4. 할인현금흐름수지표는 회계기간을 고려하지 않고, 기준시점부터 1년 단위로 작성하되, 현재가치율은 백분율 기준 소수점 첫째자리까지 계산하고 표기할 것 \par
5. 각 점포에 대해 체결된 임대계약은 시장임대료 및 해당 점포 소재지 판매시설 부동산 시장 상황을 적절히 반영하고 있음 \par
6. A자산운용사는 대형할인점 운영기업의 경우 최근 소셜커머스 기업의 대형화로 인해 성장이 둔화된 상황이므로, 향후 건물가치가 하락할 수 있다는 점을 고려하여 부동산 매입을 결정하였음. 지역별 할인율을 결정할 때, 건물가치의 운용기간 중 회수율을 고려하여 자산별 적용 할인율의 균형을 검토하여야 함 \par
7. 상기 제시된 모든 수익률은 감가상각비를 비용으로 고려하지 않은 상각전 수익에 대한 원본가치 대비 수익률임 \par
8. 해당 펀드는 대상 점포별로 다른 금융기관의 대출을 이용하고자 하는바, 각 점포별 담보가치의 균형에 유의하여야 함 \par
9. 수익률 산정 및 수익환원법 시산가액 산정에 있어 세금효과는 배제함 \par
10. 물음 2)의 할인현금수지표에는 기간별 순영업소득, 재매도 가액에서 재매도 비용을 공제한 순재매도 가액, 이자지급전 현금흐름, 현재가치율, 할인현금흐름이 포함되어야 함 \par
11. 평가개요 작성은 생략할 것 \par
\end{longquestion}
